% =============================================================================
\section{Motivation and the use of AI}
\label{sec:history}

This paper did not begin in loop quantum gravity. Since late June 2026 I had been
building an attention layer, which I call the transporter, in which three tokens interact
through one trilinear form, so that its basic cell is a triangle, and whose edges carry
operators on the Albert algebra, the exceptional Jordan algebra of $3\times3$ octonionic
Hermitian matrices. On that algebra the tensor square of a traceless element splits into
five irreducible channels, $26\otimes26=1\oplus26\oplus52\oplus273\oplus324$, and the layer
is built from them~\cite{FitchPairChannels}. In October I asked whether the next simplex,
the tetrahedron, has any place in this setting. It does. Four primitive idempotents of the
Albert algebra with equal pairwise overlaps $\frac13$ form a regular tetrahedron of pure
states that generates the whole algebra. And a quantum tetrahedron in the sense of Barbieri
and of Baez and Barrett~\cite{Barbieri1998,BaezBarrett1999}, with the exceptional group
$F_4$ in place of $\mathrm{SU}(2)$ and the $26$ on each of its four faces, has a
five-dimensional space of shapes, and the five channels label a basis of it. Its
orientation operator is the commutator of two pair Casimirs, which for $\mathrm{SU}(2)$
faces is the triple product $J_1\cdot(J_2\times J_3)$, the four-valent volume operator of
loop quantum gravity.

The volume operator is what took me to the dynamics. In a Claude session on 17 and
18~October I asked why more physicists had not talked about loop quantum gravity, and the
answer named the unsolved dynamics, the quantisation of the Hamiltonian constraint and the
anomaly-free closure of its algebra, as the central problem. In the same session I wrote
a research prompt for the rank-three Jordan family, and one of its four tasks was to define the
analogue of the Baez--Barrett volume operator and compute its spectrum, calibrated against
the published volume spectra of Brunnemann and Thiemann and of Bianchi and
Haggard~\cite{BT2006,BH2011}. In the dynamics the volume has a role of its own: in
Thiemann's construction the extrinsic curvature is obtained from the commutator of the
Hamiltonian constraint with the volume~\cite{ThiemannQSD}. The first lemma proved in the
work that followed was that this commutator vanishes on the off-shell space that Varadarajan
built for the Euclidean theory~\cite{Varadarajan2022}, and the question became whether a
Lorentzian Hamiltonian constraint can be made to close there. I did not settle it. But the
same triple product kept reappearing inside the calculation, and the Ashtekar--Lewandowski
volume entered it through the sign factors of the edge tangents that are the subject of
\cref{sec:normalform}. The questions I could settle turned out to be about the volume
operator itself. Looking for a named open problem on which these tools could act, I found
Bianchi and Haggard's question whether the volume has a gap for all spins~\cite{BH2011}.
At five-valent vertices Brunnemann and Rideout had seen gaps both open and
close~\cite{BR2008a}, and this paper explains why. The last paragraph of \cref{sec:open}
returns to dynamics, where the work began. The companion paper~\cite{FitchLimitCircle}
comes from the same calculation, continued to $\mathrm{SU}(1,1)$. \Cref{app:provenance}
describes how the results were reached.

\paragraph{Use of AI.} I posed the questions, chose the problems and the order in which to
pursue them, and set the standard a result had to meet. I directed the computations and
decided what this paper claims. Large language models were used throughout. Anthropic's
Claude was used in two ways. In sessions on claude.ai it answered my first questions about
tetrahedra in Jordan algebras, carried out the exact calculations on the rank-three Jordan
family, and began the calculation on the dynamics. Through Claude Code it wrote and ran
nearly all of the code and the certificate scripts for this paper, ran the agent campaigns,
drafted and revised the text, and gave the first written proofs of most of the results
here, among them the normal form of \cref{thm:nf-intro}, the gap closing of
\cref{thm:closing}, the sharp four-valent law of \cref{thm:sharp4} with its Stark
factorisation, and the interval certificate behind \cref{thm:z32}. Separate Claude agents,
working from the written results, refereed them one by one, and an independent Claude
referee read the whole paper. OpenAI's GPT models were used through ChatGPT and Codex.
ChatGPT produced the first document on the Albert tetrahedron, and GPT ran long
exact-arithmetic computations on the Jordan family for two further reports. Codex sessions
continued part of the dynamics calculation and audited a note on the equal-weight case.
Google's Gemini was used only as a reviewer of this paper, as were further sessions of GPT
and Claude. Review comments were adopted only after they had been checked against the
source, the rendered pages and our own code. The claims rest on the following, and not on
any model's assertion: the written proofs in this paper; the exact rational, outward-rounded
interval and rigorously bounded floating-point certificates of \cref{sec:computer}, each
shipped with its script and checksum; our recomputation of the published numbers of
Brunnemann and Rideout and of Bianchi and Haggard; and cited results. Responsibility for
the content, and for any error that remains, is mine.
